\section*{第 \ref{ch:g2:what-animals-eat} 章——动物吃什么}

\begin{solution}{exo:g2:what-animals-eat:1}
\omterm{def:g2:what-animals-eat:kinds}{植食性动物}；\omterm{def:g2:what-animals-eat:kinds}{肉食性动物}；\omterm{def:g2:what-animals-eat:kinds}{杂食性动物}。
\end{solution}

\begin{solution}{exo:g2:what-animals-eat:2}
(a) \omterm{def:g2:what-animals-eat:kinds}{植食性动物}。(b) \omterm{def:g2:what-animals-eat:kinds}{肉食性动物}。(c) \omterm{def:g2:what-animals-eat:kinds}{杂食性动物}。(d) \omterm{def:g2:what-animals-eat:kinds}{植食性动物}。(e)
\omterm{def:g2:what-animals-eat:kinds}{杂食性动物}。
\end{solution}

\begin{solution}{exo:g2:what-animals-eat:3}
扁平的磨牙——草这类\omterm{def:g1:plants-around-us:plant}{植物}很坚韧，必须磨很久
才能吞咽；也没有猎物可抓。
\end{solution}

\begin{solution}{exo:g2:what-animals-eat:4}
短而粗：咬碎\omterm{def:g2:from-seed-to-plant:seed}{种子}。弯而带钩：撕裂肉类。
细长：从树皮和缝隙中挑出昆虫。
\end{solution}

\begin{solution}{exo:g2:what-animals-eat:5}
微小的\omterm{def:g1:plants-around-us:plant}{植物}虫（蚜虫）——瓢虫吃掉大量
蚜虫，因此园丁很喜欢它。
\end{solution}

\begin{solution}{exo:g2:what-animals-eat:6}
草每口提供的营养很少，所以奶牛必须
吃上几个小时才够；肉的营养丰富，一顿饱餐
就足够狐狸享用很久。
\end{solution}

\begin{solution}{exo:g2:what-animals-eat:7}
夏天吃蚯蚓和昆虫，冬天吃浆果，因为
地面太硬，蚯蚓不好找。食谱广泛，\omterm{def:g2:what-animals-eat:kinds}{杂食性动物}
能在一种食物吃完时换吃别的——这在艰难的季节里大有帮助。
\end{solution}

\begin{solution}{exo:g2:what-animals-eat:8}
开放性答案。通常是：麻雀和雀类（喙短而粗）
吃\omterm{def:g2:from-seed-to-plant:seed}{种子}；山雀既捉昆虫也吃\omterm{def:g2:from-seed-to-plant:seed}{种子}；乌鸫在
地里寻找蚯蚓，并吃浆果。“喙与食谱的匹配”这一命题
在这里成立得很好。
\end{solution}

\begin{solution}{exo:g2:what-animals-eat:9}
\omterm{def:g2:what-animals-eat:kinds}{肉食性动物}的食谱：其他\omterm{def:g1:animals-around-us:animal}{动物}。第 1 和第 2 步——牙齿
（尖利的抓握工具）和捕猎工具（爪子、向前的眼睛）；
第 4 步已经给出了它的名称。
\end{solution}

\begin{solution}{exo:g2:what-animals-eat:10}
狐狸吃兔子，兔子吃草。没有草，
就没有兔子；没有兔子，狐狸就会挨饿。狐狸依赖草，
\emph{通过}其猎物——自然界的进食关系形成食物链。
\end{solution}
